\section{Digital Card Game} \label{app:dcg}
\subsection{Dataset Details}
\vpara{Construction Details.}
We use Aquawar framework as the basis for our interactive system. The first type of interaction is the \textbf{action} phase, where the model needs to select the fish it wants to act with and then choose the target for skill. To ensure the validity of model operations, we perform checks for valid actions. The second type of interaction is the \textbf{guess} phase, where we provide the model with known information, including fish species and skill descriptions, enemy's targets. We have two naive strategies (random and greedy search) for testing purposes. The following is a detailed definition and description of the game process.

\begin{itemize}[leftmargin=1.5em,itemsep=0pt,parsep=0.2em,topsep=0.0em,partopsep=0.0em]
    \item \textbf{Player and Cards.} It is a two-player battle game with four pet fishes (i.e., cards) in each team. The card pool consists of ten fish (Appendix~\ref{app:fish_attribute}), and both players choose four definite fish to use before the start of the game.
    \item \textbf{Initial State.} Each fish has 400 initial health, 200 initial attack power, active ability, and passive ability.
    \item \textbf{Basic Rule.} Players choose a live fish to use its active skill or normal attack on an enemy fish each round. All alive fish's passive ability will automatically trigger when meeting certain conditions.
    \item \textbf{Assertion Mechanism.} The identity of a player's fish is initially hidden. The counter-player can guess one of the player's fish's identities each round. If the counter-player guesses correctly, the player's fish's identity is revealed, and all its fish will get damaged.
    \item \textbf{Round Process.} Within a round of the game, the player for that round will first assert the identity of one opponent's fish that are alive and whose identities have not been revealed. If the assertion is correct, all of the opponent's fish that remain alive get damaged. Subsequently, the player for that round can command one alive fish to execute a normal attack or an active ability. Following this, any fish that meet the condition will unleash its passive ability.
    \item \textbf{Victory Condition.} The victory condition is to have more fish alive at the end of the game.
\end{itemize}
To balance agent engagement and game complexity simultaneously, we designed two stages of game logic. We remove the assertions in the first stage while keeping assertions in the second stage. We test all the models on both the first and second stages separately and choose the average performance for final score.

We choose two naive playing strategies as the baselines.
\begin{itemize}[leftmargin=1.5em,itemsep=0pt,parsep=0.2em,topsep=0.0em,partopsep=0.0em]
\item The first strategy is a simply random action from all available action spaces.
\item The second strategy will try to use AOE attack if possible, and continuously evaluating whether a one-hit kill is possible. Then, it attempts to use active skills and, finally, resorts to normal attacks. Overall, this strategy follows a certain pattern but may not necessarily be the most optimal one. 
\end{itemize}

\vpara{Evaluation Setup.}
For each time of the game playing, we evaluate with the following steps:

\begin{itemize}[leftmargin=1.5em,itemsep=0pt,parsep=0.2em,topsep=0.0em,partopsep=0.0em]
    \item \textbf{Initialization.} We initiated the modified game logic environment, which uses pybind to compile, and the baseline game agent under the Ubuntu 20.04 environment.
    \item \textbf{Interaction.} We place rule descriptions in the instruction prompt according to different game stages, and the LLM agent interacts and competes strategically with the baseline within the game logic environment. We give the LLM agent five chances to respond in the correct format. It will be immediately deemed defeated if it fails to output legal actions within the given number of attempts. At the same time, we encourage the model to output its reasoning process in CoT.
    \item \textbf{Result Calculation.} During the Interaction process, we will record the entire game process for battle playback and calculate the game results to obtain the metrics for the task.
\end{itemize}

\vpara{Metrics.}
Our comprehensive evaluation uses metrics that range from basic gameplay elements such as the wining rounds \textbf{(Win Round)} , total played rounds \textbf{(Total Round)}, winning rate \textbf{(Win Rate)} , the total damage inflicted compared to total health \textbf{(Damage Rate)}, 
and ultimately we provide a final reward score according to the above metrics: $$\mathbf{reward} = 0.7 \times \mathbf{metric}_{winrate} + 0.3 \times \mathbf{metric}_{damagerate}$$

\subsection{The Attributes of Fish} \label{app:fish_attribute}
The game has ten kinds of fish according to the game rules.
\begin{itemize}[leftmargin=1.5em,itemsep=0pt,parsep=0.2em,topsep=0.0em,partopsep=0.0em]
\item \textbf{Spray} 
    
    - \textbf{Counter (Passive)}: Inflicts 30 damage to the attacker when a teammate's health is below 30\%
    
    - \textbf{AOE (Active)}: Attacks all enemies for 35\% of its attack points.

\item \textbf{Flame}

    - \textbf{Counter (Passive)}: Inflicts 30 damage to the attacker when a teammate's health is below 30\%
    
    - \textbf{Infight (Active)}: Inflicts 75 damage on one living teammate and increases your attack points by 140.

\item \textbf{Eel}

   - \textbf{Deflect (Passive)}: Distributes 70\% damage to teammates and takes 30\% when attacked. Gains 40 attack points after taking 200 damage accumulated.
    
   - \textbf{AOE (Active)}: Attacks all enemies for 35\% of its attack points.

\item \textbf{Sunfish}

    - \textbf{Deflect (Passive)}: Distributes 70\% damage to teammates and takes 30\% when attacked. Gains 40 attack points after taking 200 damage accumulated.
    
    - \textbf{Infight (Active)}: Inflicts 75 damage on one living teammate and increases your attack points by 140.

\item \textbf{Barracuda}

    - \textbf{Reduce (Passive)}: There is a 30\% chance to avoid any incoming damage each time.
    
    - \textbf{Crit (Active)}: Deals 120 CRITICAL damage to an enemy.

\item \textbf{Mobula}

    - \textbf{Reduce (Passive)}: There is a 30\% chance to avoid any incoming damage each time.
    
    - \textbf{Subtle (Active)}: Choose a teammate or yourself to reduce the damage taken by 70\% when attacked, and increase its attack points by 20.

\item \textbf{Octopus}

    - \textbf{Heal (Passive)}: Regain 20 health points if the health is still greater than 0 when attacked. 
    
    - \textbf{Infight (Active)}: Inflicts 75 damage on one living teammate and increases your attack points by 140.

\item \textbf{Whiteshark}

    - \textbf{Heal (Passive)}: Regain 20 health points if the health is still greater than 0 when attacked.
    
    - \textbf{Crit (Active)}: Deal 120\% CRITICAL damage of your attack power to the enemy with the lowest health. If the target's health is below 160, increase the CRITICAL damage to 140\%. 

\item \textbf{Hammerhead}

    - \textbf{Explode (Passive)}: Deal 40 damage to the source when attacked but not died. When the health is below 20\%, increase its attack points by 15. 
    
    -  \textbf{Crit (Active)}: Deal 120\% CRITICAL damage of your attack power to the enemy with the lowest health. If the target's health is below 160, increase the CRITICAL damage to 140\%. 

\end{itemize}
As can be seen, there is overlap among the active and passive skills of different pet fish, which is done to better conceal the identity information of pet fish in the game and increase the strategic aspects of the game.

\subsection{Prompt Example}
We use the following format of prompts for actions:
\lstset{
    backgroundcolor=\color[RGB]{245,245,244},
    breaklines=true,
    basicstyle=\ttfamily\small
}\begin{lstlisting}
This is a two-player battle game with four pet fish on each team. The types of fish may vary.
Each fish has its 400 initial health, 200 attack power, active ability, and passive ability.
You can choose a live fish to use its active skill or normal attack (causing half of attack power as damage) on an enemy fish each round.
When the conditions are met, the fish's passive ability will automatically trigger, regardless of whether it is chosen.
Your fish's identity is initially hidden. The enemy can guess one of your fish's identity in each round. If the enemy guesses right, your fish's identity is revealed, and each of your fish will get 50 damage.
The victory condition is to have more fish alive at the end of the game.

The following are the four types of your pet fish:
{'spray': {'passive': "Counter: Deal 30 damage to attacker when a teammate's health is below 30%. ", 'active': 'AOE: Attack all enemies for 35% of its attack points.'}, 'flame': {'passive': "Counter: Deal 30 damage to attacker when a teammate's health is below 30%. ", 'active': "Infight: Attack one alive teammate for 75 damage and increases your attack points by 140. Notice! You can't attack yourself or dead teamate! "}, 'eel': {'passive': 'Deflect: Distribute 70% damage to teammates and takes 30% when attacked. Gains 40 attack points after taking 200 damage accumulated. ', 'active': 'AOE: Attack all enemies for 35% of your attack points.'}, 'sunfish': {'passive': 'Deflect: Distribute 70% damage to teammates and takes 30% when attacked. Gains 40 attack points after taking 200 damage accumulated. ', 'active': "Infight: Attack one alive teammate for 75 damage and increases your attack points by 140. Notice! You can't attack yourself or dead teamate! "}}

The following are the four types of enemy's pet fish:
{'spray': {'passive': "Counter: Deal 30 damage to attacker when a teammate's health is below 30%. ", 'active': 'AOE: Attack all enemies for 35% of its attack points.'}, 'flame': {'passive': "Counter: Deal 30 damage to attacker when a teammate's health is below 30%. ", 'active': "Infight: Attack one alive teammate for 75 damage and increases your attack points by 140. Notice! You can't attack yourself or dead teamate! "}, 'eel': {'passive': 'Deflect: Distribute 70% damage to teammates and takes 30% when attacked. Gains 40 attack points after taking 200 damage accumulated. ', 'active': 'AOE: Attack all enemies for 35% of your attack points.'}, 'sunfish': {'passive': 'Deflect: Distribute 70% damage to teammates and takes 30% when attacked. Gains 40 attack points after taking 200 damage accumulated. ', 'active': "Infight: Attack one alive teammate for 75 damage and increases your attack points by 140. Notice! You can't attack yourself or dead teamate! "}}

Play the game with me. In each round, you should output your thinking process, and return your move with following JSON format:
{'pick_fish': 'pick an alive fish, you should give the name of the alive fish', 'action': 'choose from [normal, active]', 'target_position': "target's position, you must choose from [0,3]"}

Notice! You must return your move in each round. Otherwise, you will be considered defeated.
\end{lstlisting}

We use the following format of prompts for assertions in stage2:
\lstset{
    backgroundcolor=\color[RGB]{245,245,244},
    breaklines=true,
    basicstyle=\ttfamily\small
}\begin{lstlisting}
This is a two-player battle game with four pet fish in each team. The types of fish may vary.
Each fish has its initial health, attack power, active ability, and passive ability.
All fish's identities are initially hidden. You should guess one of the enemy fish's identities in each round. If you guess right, the enemy fish's identity is revealed, and each of the enemy's fish will get 50 damage. You can only guess the identity of the live fish.
The victory condition is to have more fish alive at the end of the game.

The following are the four types of your pet fish:
{'spray': {'passive': "Counter: Deal 30 damage to attacker when a teammate's health is below 30%. ", 'active': 'AOE: Attack all enemies for 35% of its attack points.'}, 'flame': {'passive': "Counter: Deal 30 damage to attacker when a teammate's health is below 30%. ", 'active': "Infight: Attack one alive teammate for 75 damage and increases your attack points by 140. Notice! You can't attack yourself or dead teamate! "}, 'eel': {'passive': 'Deflect: Distribute 70% damage to teammates and takes 30% when attacked. Gains 40 attack points after taking 200 damage accumulated. ', 'active': 'AOE: Attack all enemies for 35% of your attack points.'}, 'sunfish': {'passive': 'Deflect: Distribute 70% damage to teammates and takes 30% when attacked. Gains 40 attack points after taking 200 damage accumulated. ', 'active': "Infight: Attack one alive teammate for 75 damage and increases your attack points by 140. Notice! You can't attack yourself or dead teamate! "}}

The following are the four types of enemy's pet fish:
{'spray': {'passive': "Counter: Deal 30 damage to attacker when a teammate's health is below 30%. ", 'active': 'AOE: Attack all enemies for 35% of its attack points.'}, 'flame': {'passive': "Counter: Deal 30 damage to attacker when a teammate's health is below 30%. ", 'active': "Infight: Attack one alive teammate for 75 damage and increases your attack points by 140. Notice! You can't attack yourself or dead teamate! "}, 'eel': {'passive': 'Deflect: Distribute 70% damage to teammates and takes 30% when attacked. Gains 40 attack points after taking 200 damage accumulated. ', 'active': 'AOE: Attack all enemies for 35% of your attack points.'}, 'sunfish': {'passive': 'Deflect: Distribute 70% damage to teammates and takes 30% when attacked. Gains 40 attack points after taking 200 damage accumulated. ', 'active': "Infight: Attack one alive teammate for 75 damage and increases your attack points by 140. Notice! You can't attack yourself or dead teamate! "}}

Play the game with me. In each round, you should output your thinking process, and return your move with following JSON format:
{'guess_type': "the enemy's fish type you may guess", 'target_position': "guess target's position, you must choose from [0,3]"}

Notice! You must return your move in each round. Otherwise, you will be considered defeated.
\end{lstlisting}


\subsection{Battle Generation}

% \todo{Newly added. Plz check @xiao}

In the main text of our study, we fixed certain presets for fair evaluation, allowing all models to be tested using the same battle scenarios. However, the generation of these tasks can indeed incorporate randomness. The randomness in a game primarily manifests in the following aspects:



\begin{itemize}
    \item \textbf{Team Composition}: The selection of pet fish for both sides, including their number and types.
    \item \textbf{Attribute Determination}: The specific values for each pet fish’s attributes on both sides.
\end{itemize}

We allow the selection of different difficulty levels to generate random battles. In this section, we will explain how this concept of difficulty works.

\subsubsection{Combat Power}

To effectively assess combat power, it's crucial to recognize its fundamental principle: in a confrontation between two teams, the one with a superior combat power typically has a higher probability of emerging victorious.

Let's denote the health points of a fish \( c \) as \( HP_{c} \) and its attack power (the expected damage of a standard attack) as \( ATK_{c} \).

Imagine a straightforward scenario where two fish are pitted against each other: one from our team, denoted as \( f \) (friendly), and the other from the opposing team, denoted as \( h \) (hostile). They engage in combat, dealing damage to each other simultaneously. The duration each can withstand the battle is represented by \( \frac{HP_{f}}{ATK_{h}} \) and \( \frac{HP_{h}}{ATK_{f}} \), respectively. In this context, our team is more likely to triumph if \( \frac{HP_{f}}{ATK_{h}} \ge \frac{HP_{h}}{ATK_{f}} \), or equivalently, \( HP_{f} \cdot ATK_{f} \ge HP_{h} \cdot ATK_{h} \).

Hence, for a single-fish team, we define its combat power as \( HP \cdot ATK \), aligning with the aforementioned criteria.

When considering a team \( T \) comprising multiple fish, and excluding any special abilities, each combat turn involves choosing one of our fish to attack an opponent's fish. This allows us to treat the team collectively, with the total health being the cumulative health of all fish and the average attack power being considered (under normal circumstances, the highest attack power would be the logical choice, but given that opponents often target the fish with the highest attack power first, reducing its longevity, we assume an equal attack frequency across all fish).

This leads us to the definition of combat power \( Power(T) \) for a team \( T \):

\[ Power(T) = \frac{1}{\#T} \sum_{c\in T}HP_c \sum_{c\in T}ATK_c \]

\subsubsection{Difficulty: A Ratio of Combat Powers}

The difficulty of a game is quantified by the ratio of combat powers between the opposing teams:

\[ \rho(H | F) = \frac{Power(H)}{Power(F)} \]

In this equation, \( F \) symbolizes the friendly team, while \( H \) represents the hostile team. Notably, a ratio of 1 signifies a balanced or normal difficulty level, indicating parity between the teams.

Furthermore, this difficulty metric implies that \( \rho(T_1 | T_2) \) instances of team \( T_2 \) would be equivalent in power to team \( T_1 \).

With this framework, researchers can accurately set specific difficulty levels for evaluating the win rates of various strategies, thus effectively measuring their efficacy.
